\documentclass[11pt,a4paper]{article}

% ==== Fonts & CJK ====
\usepackage{xeCJK}
\setCJKmainfont{Songti SC}
\setCJKsansfont{PingFang SC}
\setCJKmonofont{Menlo}
\setmainfont{Times New Roman}
\setmonofont{Menlo}

\usepackage[margin=2.2cm]{geometry}
\usepackage{amsmath, amssymb, amsthm, mathtools}
\usepackage{bm}
\usepackage{booktabs, tabularx, array, multirow}
\usepackage{graphicx}
\usepackage{xcolor}
\usepackage{hyperref}
\usepackage{enumitem}
\usepackage{caption}
\usepackage{fancyhdr}
\usepackage{titling}
\usepackage{tcolorbox}
\usepackage{tikz}
\usetikzlibrary{arrows.meta, positioning, shapes, fit}

\hypersetup{
  colorlinks=true,
  linkcolor=blue!60!black,
  urlcolor=blue!60!black,
  citecolor=blue!60!black,
  pdftitle={RS+BCH 级联码低时延译码课题解读},
  pdfauthor={陈诗阳}
}

\pagestyle{fancy}
\fancyhf{}
\rhead{\footnotesize RS+BCH 级联码课题解读}
\lhead{\footnotesize 基于 IEEE TIT 2025 LLOSD}
\cfoot{\thepage}

\renewcommand{\baselinestretch}{1.2}

\definecolor{accent}{RGB}{30,80,160}
\definecolor{accent2}{RGB}{160,40,40}
\definecolor{accent3}{RGB}{40,120,60}

\newtcolorbox{keybox}[1][]{colback=blue!5, colframe=accent, boxrule=0.4mm,
  arc=1mm, left=2mm, right=2mm, top=1mm, bottom=1mm, #1}
\newtcolorbox{riskbox}[1][]{colback=red!5, colframe=accent2, boxrule=0.4mm,
  arc=1mm, left=2mm, right=2mm, top=1mm, bottom=1mm, #1}
\newtcolorbox{actionbox}[1][]{colback=green!5, colframe=accent3, boxrule=0.4mm,
  arc=1mm, left=2mm, right=2mm, top=1mm, bottom=1mm, #1}

\begin{document}

% ============ Cover ============
\begin{center}
{\Large\bfseries 导师 RS+BCH 级联码低时延译码课题解读}\\[4pt]
{\large 技术方案、待对齐问题与执行路线图}\\[16pt]
{\normalsize 陈诗阳\\
基于 IEEE TIT 2025 论文《Efficient Ordered Statistics Decoding of BCH Codes Without Gaussian Elimination》\\
\vspace{2pt}
\texttt{https://github.com/a-yeyang/efficient-osd-bch-no-ge}}\\[10pt]
\rule{0.9\textwidth}{0.4pt}\\[6pt]
{\footnotesize 撰写于 2026-07-23 · 基于 Claude 4.7 Opus + Claude Code CLI}\\
\rule{0.9\textwidth}{0.4pt}
\end{center}

\vspace{6pt}
\begin{keybox}[title=\textbf{一句话总结导师要做的事}]
构造一个 \textbf{RS + BCH 级联译码器}，内码用 LLOSD（Lagrange 版 OSD），
外码用软/硬判决 RS，让两级 Lagrange 插值\textbf{共享代数计算}，
从而在加入 BCH 内码后，级联链路的译码时延相比纯 RS 只增加
\textbf{$\leq$10\%}。
\end{keybox}

% ============ Section 1 ============
\section{导师要求逐字翻译}

先把导师的手写要求转成精确的技术含义。

\begin{table}[h!]
\centering
\begin{tabularx}{\textwidth}{lX}
\toprule
\textbf{导师原文} & \textbf{精确技术含义} \\
\midrule
RS+BCH 级联码 &
外码 RS (定义在 $\mathbb{F}_{2^m}$)，内码 BCH (定义在 $\mathbb{F}_2$)。
发端: BCH 编码 $\to$ RS 编码；收端: 反向 \\
开销 12\% (码率 0.88) &
$R_{\text{total}} = R_{\text{RS}} \cdot R_{\text{BCH}} \approx 0.88$ \\
码长 128 和 255 &
两种典型 config（对应 $\text{GF}(2^7)$ 和 $\text{GF}(2^8)$），覆盖光通信 OTU / 短码场景 \\
AWGN + PAM4 &
4 电平幅度调制（$-3, -1, +1, +3$ Gray 编码），每符号 2 bit；
用于 400G/800G 光模块 \\
RS 硬盘和软判 & RS 两条路径: BM (硬判) + KV/LCC-BR (软判) \\
BCH 用 LLOSD & 直接用 IEEE TIT 2025 论文里的算法 \\
核心思想 &
利用 BCH $\subset$ RS 的代数嵌入，用 Lagrange 替 GE，级联外码/内码共享插值计算 \\
BCH 时延增加 $\leq$10\% &
硬性 KPI: $(T_{\text{RS+BCH}} - T_{\text{RS only}}) / T_{\text{RS only}} \leq 10\%$ \\
导出 Matlab 程序 & 最终代码是 Matlab (通信仿真行业惯例) \\
\bottomrule
\end{tabularx}
\caption{导师需求逐条翻译}
\end{table}

% ============ Section 2 ============
\section{导师 idea 到底是什么}

\subsection{表面 idea (字面读)}

字面上，导师说``把 Lagrange 那套用来构造级联码外码的 RS 系统生成矩阵''。
但这句话有点绕: 外码本身就是 RS，它本来就用 Lagrange 或者移位寄存器构造，
跟论文没关系。

\subsection{真正的 idea (技术核心)}

论文 LLOSD 的核心洞察是: BCH $\subset$ RS，所以\textbf{BCH 可以用 RS 的代数
(Lagrange 插值) 来译码}。导师想把这个洞察\textbf{推广到级联码}:

\begin{keybox}[title=\textbf{核心创新点}]
\textbf{级联码译码时，内外码之间共享 Lagrange 插值的中间计算结果}。

因为内码 BCH 的 LLOSD 用 Lagrange (构造 $\mathbf{G}_{\text{RS,inner}}$)，
外码 RS 的 LCC-BR/KV 也用 Lagrange，这两次 Lagrange 插值\textbf{共用同一个
代数结构} (同一个 GF($2^m$) 上的 $\alpha$ 幂序列)，
大量中间计算 (locators $\alpha^j$, denominator products $\prod(\alpha^i - \alpha^j)$,
Lagrange 基函数 $T_j(\alpha^i)$)\textbf{可以复用}。
\end{keybox}

这正是论文 §V 里 HSD 算法 (Algorithm 2) 的思想: LLOSD + LCC-BR 集成时，
因为两者都是 RS 域算法，共享代数计算。\textbf{导师是要把 HSD 从
``BCH 单码内部''推广到``级联码内外双码''}。

\subsection{译码链路示意}

\begin{center}
\begin{tikzpicture}[node distance=0.7cm, >=Stealth]

\node[draw, rounded corners, minimum width=3cm, minimum height=0.9cm, fill=blue!5]
     (llr) {接收端 PAM4 $\to$ bit-LLR};

\node[draw, rounded corners, minimum width=5.2cm, minimum height=1.1cm,
     fill=blue!10, below=of llr, align=center]
     (bch) {\textbf{内码 BCH 译码 (LLOSD)}\\
             {\footnotesize $\mathbf{G}_{\text{RS,inner}}$ 由 Lagrange 插值构造}};

\node[draw, rounded corners, minimum width=5.2cm, minimum height=1.1cm,
     fill=blue!10, below=of bch, align=center]
     (rs) {\textbf{外码 RS 译码 (BM 或 LCC-BR)}\\
             {\footnotesize 也需要 Lagrange 插值 \& 相同的 $\alpha^j$ 表}};

\node[draw, rounded corners, minimum width=3cm, minimum height=0.9cm,
     fill=green!10, below=of rs]
     (out) {恢复的消息};

\node[draw, dashed, rounded corners, fit=(bch)(rs),
     inner sep=0.3cm, label={right:{\textbf{Lagrange 共享区}}}] {};

\draw[->] (llr) -- (bch);
\draw[->] (bch) -- (rs);
\draw[->] (rs) -- (out);

\end{tikzpicture}
\end{center}

% ============ Section 3 ============
\clearpage
\section{具体要做的 6 件事 (拆解到可执行)}

\begin{table}[!ht]
\centering
\begin{tabularx}{\textwidth}{clX}
\toprule
\# & \textbf{任务} & \textbf{关键细节} \\
\midrule
1 & 搭建 RS+BCH 级联码框架 &
选参数满足 $R_{\text{total}} = R_{\text{RS}} \cdot R_{\text{BCH}} \approx 0.88$。
{\bf n=255 候选:} RS(255, 239) + BCH(255, 239)，总码率 0.879。
{\bf n=128 候选:} 需要在 GF($2^7$) 上试。
先 BCH 后 RS。每 $m=8$ bit 打成一个 RS 符号 \\
\midrule
2 & PAM4 调制 + AWGN 信道 &
PAM4 电平 $\{-3, -1, +1, +3\}$ Gray 编码，每符号 2 bit。
接收端要计算 2 个 bit-LLR (不是简单高斯 LLR，需考虑 4 电平混叠) \\
\midrule
3 & 内码 BCH 译码 (LLOSD) &
把复现的 LLOSD 代码 port 到 Matlab。
输入 $n$ 个 bit-LLR，输出 $n$ 个硬比特 + 每个 RS 符号的可靠性度量 \\
\midrule
4 & 外码 RS 译码 (两版) &
{\bf 硬判:} BM + Chien + Forney。
{\bf 软判:} KV (Kötter–Vardy 2003) 或 LCC-BR (Xing–Chen–Bossert 2020)，
后者更快、硬件友好 \\
\midrule
5 & \bf Lagrange 共享 (核心创新) &
在 LLOSD 里已经算了 $T_j(\alpha^i)$ 表，
外码 RS 的 KV/LCC-BR 也需要类似的表。
实现\textbf{一次计算、两次使用}的 Lagrange 缓存层。
量化时延节省 \\
\midrule
6 & 仿真两个 KPI &
{\bf KPI 1 (BER-SNR):} 纯 RS 硬判 / 纯 RS 软判 / RS+BCH 级联 (RS硬+BCH软) /
RS+BCH 级联 (RS软+BCH软) 四条曲线。预期级联码有 1--2 dB 编码增益。
{\bf KPI 2 (时延):} 证明 $\Delta T / T_{\text{RS only}} \leq 10\%$ \\
\bottomrule
\end{tabularx}
\end{table}

% ============ Section 4 ============
\section{课题的科研价值}

从信息论/通信原理视角看，导师的选题非常有品味:

\begin{itemize}[leftmargin=1.5em]
\item \textbf{场景明确}: 400G/800G 光模块用 PAM4 + RS(544, 514) KP4 已是 IEEE 802.3
  标准，但下一代 200G/lane 需要更强的 FEC。\textbf{在标准 KP4 RS 外面套一层 BCH 内码}
  是学术界+工业界都在探索的方向 (OIF 400ZR+, 800ZR)。
\item \textbf{痛点真实}: 级联码的\textbf{时延}是决定它能不能上标准的第一障碍。
  KP4 RS 时延已经很紧 (几十 ns)，任何 inner code 都必须近乎零时延增加。
  10\% 是硬 KPI。
\item \textbf{切入点犀利}: BCH 内码软判可以带来 1--2 dB 增益，
  但传统 OSD 时延爆炸。\textbf{LLOSD 恰好把 BCH 软判时延压到与硬判同数量级}
  --- 这正是导师看中这篇论文的原因。
\item \textbf{可推广性}: 若证明 Lagrange 共享真能压到 10\% 以内，
  可以直接推广到 RS+RS (RS$^2$, 800ZR+) / LDPC+BCH / Polar+CRC 级联。
\end{itemize}

% ============ Section 5 ============
\section{隐藏的技术风险}

\begin{riskbox}[title=\textbf{Lagrange 共享的可行性风险}]
LLOSD 里的 $\mathbf{G}_{\text{RS}}$ 是\textbf{每次译码基于当前 MRIP 重新构造}的
(因为 MRIP 依赖 LLR)。这意味着\textbf{内外码之间的 Lagrange 共享并不是
``预计算一次''，而是每次都要 recompute}。

要达到 10\% 时延增加目标，可能需要:
\begin{itemize}[leftmargin=1em, itemsep=1pt, topsep=2pt]
\item 只共享 $\alpha$ 幂表 ($2^m$ 个元素，可完全预计算)
\item 只共享 denominator products 的公共部分
\item \textbf{内外码用同一个 primitive polynomial}
  (这是能不能共享的先决条件！)
\item 用 look-up table 替换乘法 (对 $m \leq 8$ 可行)
\end{itemize}
\end{riskbox}

\begin{riskbox}[title=\textbf{PAM4 LLR 的正确性风险}]
PAM4 的 LLR 不能用 BPSK 的 $2y/\sigma^2$ 公式。
每个 PAM4 符号对应两个 bit，两个 bit 的 LLR 都要 marginalize 掉另一个 bit:
\[
\text{LLR}(b_1 | y) = \log \frac{\sum_{s: b_1(s)=0} \exp(-\|y - s\|^2 / 2\sigma^2)}
                                {\sum_{s: b_1(s)=1} \exp(-\|y - s\|^2 / 2\sigma^2)}
\]
Gray 编码保证相邻电平差 1 bit，但 LLR 的 max-log 近似仍要注意边界。
BCH 是 bit-level 码，LLR 输入必须是 per-bit 的。
\end{riskbox}

\begin{riskbox}[title=\textbf{参数选择的风险}]
$n \in \{128, 255\}$ 是电通信/光通信惯例。但 BCH 是 primitive narrow-sense 码，
真实的 BCH 码长是 $2^m - 1$，即 127 和 255。要么:
\begin{itemize}[leftmargin=1em, itemsep=1pt, topsep=2pt]
\item 用扩展 BCH (extended BCH, 加一位奇偶校验) 得到 $n=128$
\item 或用 shortened BCH (删除信息位)
\end{itemize}
两种做法在 LLOSD 里的具体实现细节不同，要跟导师确认。
\end{riskbox}

% ============ Section 6 ============
\section{给你的执行路线 (三阶段)}

\begin{table}[h!]
\centering
\begin{tabularx}{\textwidth}{llX}
\toprule
\textbf{阶段} & \textbf{时长} & \textbf{交付物} \\
\midrule
Phase 1 & 1--2 周 &
搭级联框架 + PAM4 信道 + 硬判 baseline (纯 RS BM + BCH BM 级联)。
BER-SNR 曲线出结果，且与文献对得上 \\
\midrule
Phase 2 & 2--3 周 &
port LLOSD 到 Matlab (Python 复现代码是现成参考) + 外码 RS 软判 (LCC-BR)。
四条 BER-SNR 曲线画出来 \\
\midrule
Phase 3 & 2--3 周 &
Lagrange 共享层实现 + 时延 profiling + 冲 10\% KPI。
若达不到，退而求其次: 分析 gap，提出改进 \\
\bottomrule
\end{tabularx}
\end{table}

% ============ Section 7 ============
\section{立刻要跟导师对齐的问题}

\begin{actionbox}[title=\textbf{Q1. 技术路线确认}]
``老师，我理解 Lagrange 共享是指内码 LLOSD 和外码 RS 软判都用 Lagrange 插值，
我们把这两次插值的公共计算提取出来。是这个意思吗？还是您有别的想法？''
\end{actionbox}

\begin{actionbox}[title=\textbf{Q2. 时延 KPI 基线}]
``BCH 时延增加 $\leq$10\% 是相对于哪个基线？
是纯 RS 硬判 BM 还是纯 RS 软判 LCC-BR？
测量口径是 clock cycles、wall-clock $\mu$s，还是 F$_{2^m}$ ops？''
\end{actionbox}

\begin{actionbox}[title=\textbf{Q3. 工作量范围}]
``两个码长 128 和 255 是并列 deliverable 还是选一个重点做？
若时间紧，是否可以先完成 $n=255$ 方案，然后 $n=128$ 作为扩展？''
\end{actionbox}

\begin{actionbox}[title=\textbf{Q4. RS 软判具体算法}]
``外码 RS 软判是用 Kötter--Vardy (2003) 还是 LCC-BR (Xing--Chen--Bossert 2020)?
后者更适合硬件、复杂度更低，但两者性能不一样。''
\end{actionbox}

\begin{actionbox}[title=\textbf{Q5. Matlab 版本 \& Toolbox}]
``可以用 Matlab Communications Toolbox 的 \texttt{comm.RSEncoder/Decoder}
和 \texttt{comm.BCHEncoder/Decoder} 吗？还是要求所有译码器都自己实现？''
\end{actionbox}

\begin{actionbox}[title=\textbf{Q6. 具体参数}]
``可以给我一组精确的 $(n_{\text{outer}}, k_{\text{outer}}, n_{\text{inner}},
k_{\text{inner}})$ 目标吗？我准备了几个候选:

\begin{itemize}[leftmargin=1em, itemsep=1pt, topsep=1pt]
\item \bf n=255 方案A: RS(255, 239) + BCH(255, 239)，总码率 0.879
\item \bf n=255 方案B: RS(255, 235) + BCH(255, 247)，总码率 0.891
\item \bf n=128 方案: RS(127, 119) + BCH(127, 120)，总码率 0.885
\end{itemize}
您有 preference 吗？''
\end{actionbox}

% ============ Section 8 ============
\section{总结}

\begin{keybox}[title=\textbf{全文核心}]
\begin{enumerate}[leftmargin=1.5em, itemsep=2pt, topsep=2pt]
\item 导师课题的\textbf{真正创新点}是：把 IEEE TIT 2025 的 LLOSD (BCH 单码内的
Lagrange 替 GE) 推广到 \textbf{RS+BCH 级联码}，让内外码共享 Lagrange 插值计算。
\item 目标 KPI: \textbf{BCH 时延增加 $\leq$10\%}，这是能不能上光通信标准的门槛。
\item 应用场景: 下一代 400G/800G 光模块 (PAM4 + KP4 RS 外套 BCH 内码)。
\item 工作量: 3 个阶段、约 6--8 周，最终交付 Matlab 代码 + BER/时延仿真报告。
\item 主要风险: Lagrange 共享的可行性、PAM4 LLR 计算、BCH 码长 128 的处理方式。
\item 下一步: \textbf{先问导师 6 个对齐问题}，不要盲动。
\end{enumerate}
\end{keybox}

\vspace{6pt}
\hrule
\vspace{4pt}
\noindent\footnotesize{
\textbf{文档说明}\\
本文档基于陈诗阳对 IEEE TIT 2025 论文的 Python 端到端复现
(\url{https://github.com/a-yeyang/efficient-osd-bch-no-ge})
以及对导师需求的技术解读整理。
所用大模型: Claude 4.7 Opus；AI 工具: Claude Code CLI。生成时间: 2026-07-23。
}

\end{document}
